\vspace{-2mm}\section{Characteristics of Reasoning Hallucination: Process vs. Result}\vspace{-1mm}
\label{sec:characteristics}

The emergence of reasoning hallucination signifies a paradigm shift in the study of LLM reliability, the locus of concern has shifted from \textit{what the model knows} to \textit{how the model thinks}.
This section contrasts reasoning hallucination with traditional hallucination to highlight its process-centric nature.

\subsection{Traditional Hallucination: Result-Centric}

Traditional hallucination, such as Factuality and Faithfulness Hallucination, is fundamentally result-centric.
Its evaluation is external and binary: a model's final answer $A$ is compared against verifiable real-world knowledge $W$ ($A \not\subset W$) or a provided source context $C$ ($A \not\subset C$).
The defect is one of \textit{knowledge}, where either the model lacks the correct fact or fails to faithfully reproduce the context.

\subsection{Reasoning Hallucination: Process-Centric}

In contrast, reasoning hallucination is fundamentally process-centric.
The primary object of scrutiny is not the correctness of the final answer $A$, but the validity and optimality of the reasoning chain $R = (r_1, \dots, r_N)$ used to derive it, which is a defect of \textit{logic}.
A model may produce a correct answer $A$ via a flawed process $R$, or a seemingly coherent process $R$ may lead to an incorrect answer~\cite{joshi2024llms,yao2025reasoning,cuesta2025lrms}.

\subsection{Core Differentiator: The Evaluability Paradox}

This shift from result to process introduces the core challenge defined in this work: the \textbf{Reasoning Evaluability Paradox}.
While traditional factual errors are at least detectable (e.g., via exact match or external search), reasoning flaws are perniciously covert.
The long, complex, and seemingly confident CoT generated by an LRM is often highly persuasive~\cite{cheng2025chain}.
This persuasiveness, a desired outcome of CoT training, ironically obscures the very flaws it generates, making detection by both human reviewers and automated tools exceptionally difficult.
The mechanism used to enhance capability simultaneously degrades our ability to assess its reliability.
